\section{Related Work}

The work most directly comparable to ours is Yim et al.~\cite{yim2023entanglement}, which detects leg entanglement from proprioception and, when a leg becomes snagged, executes a disentanglement maneuver as that leg advances through its swing phase. Their system succeeds in 14 of 16 laboratory trials spanning obstacle stiffness and geometry, and it also holds up in outdoor tests, so entanglement is a real and tractable failure mode for legged locomotion. Detection and reaction are largely rule-based and reactive, coupled to gait-phase timing rather than to a learned temporal model, and the pipeline treats sensing and control as one loop. We keep the same proprioceptive premise but change three things. First, we integrate evidence over a fixed $0.40$~s window with a learned causal model instead of applying an instantaneous threshold, which lets brief resistive transients accumulate into a decision rather than trigger one immediately. Second, we produce structured outputs: which of the four legs is affected and a calibrated per-leg severity in $[0,1]$, rather than a single binary event. Third, we treat detection as a self-contained perception problem that consumes a published state estimate, so the detector is decoupled from any actuation policy. We do not claim our detector reacts faster or acts on the world; the contrast is in what is inferred and how it is reported, not in closing a control loop.

Learning from proprioception has become standard for legged control. Hwangbo et al.~\cite{hwangbo2019learning} and Lee et al.~\cite{lee2020learning} show that policies trained largely on joint and inertial signals transfer to hardware and cope with terrain that was never seen at training time, which establishes that proprioceptive streams carry enough information to reason about contact and disturbance. A related line estimates contact state and actuator faults from the same signals: Camurri et al.~\cite{camurri2017contact} infer foot contact probabilistically for state estimation, and Bledt et al.~\cite{bledt2018contact} detect contact events to inform gait. Entanglement differs from these in kind. A snag is an externally imposed kinematic constraint, applied at an unpredictable time and place, that resists a leg's intended motion without necessarily changing its nominal contact schedule. It therefore has to be localized to a specific leg and graded by how strongly it resists, which is what motivates our per-leg attribution and physics-grounded severity index rather than a single contact-versus-no-contact readout.

For the temporal model itself we build on temporal convolutional networks~\cite{bai2018tcn} and the dilated causal convolutions introduced for audio generation in WaveNet~\cite{oord2016wavenet}. A causal, dilated design fits streaming proprioception for two reasons. Causality means the prediction at each timestep uses only past and present samples, so the model can run online without waiting for future data. Dilation stacks a modest number of layers into a receptive field wide enough to span a stride, which lets the network relate the onset of resistance to its persistence over the window without the sequential dependency of a recurrent model. These properties, rather than raw capacity, are why the architecture suits an at-every-timestep detector on a mobile robot.

A detector is only usable if its confidence score can be turned into a threshold an operator trusts, which places probability calibration on the critical path. Modern networks are frequently over-confident, and post-hoc methods such as temperature scaling~\cite{guo2017calibration} and the earlier calibration analysis of Niculescu-Mizil and Caruana~\cite{niculescu2005predicting} recalibrate scores after training so that a reported probability better reflects an empirical frequency. We adopt temperature scaling in the same spirit, mainly to de-saturate probabilities that would otherwise pin near $1.0$ and force the operating threshold to an unusable extreme; we report its effect in full, including where it does not improve every calibration measure. This lets us choose an operating point deliberately rather than by inspecting a saturated score.